\section{Applications}\label{sec:applications}
The theories reviewed here are used in four broad application areas. For each, we state which regime of the parameter space is relevant, based on the characteristic numbers derived in the previous sections.

\subsection{Laser ultrasonics and short-pulse heating}
Laser-generated ultrasound exploits the thermal stress produced by pulsed heating to generate elastic waves \citep{Scruby1990}. For nanosecond pulses on bulk metals, $\omega\tau_q\ll1$ and the classical coupled theory with surface heating is adequate for the generated ultrasound; the thermal diffusion length over the pulse duration, $\sqrt{\alpha t_p}$, sets the source depth. For femtosecond--picosecond pulses, the heat-carrier relaxation (electron--phonon coupling, in metals) occurs on time scales where hyperbolic or two-temperature models become relevant, and thin-film heating is modelled with lagging, fractional or memory-dependent laws; analyses of bi-layers under pulsed heating combine these \citep{Li2019}. The benchmark thermal shock of Section~\ref{sec:transient} shows what to look for: the stress peak at the elastic front and its sensitivity to the heat law, here about 20\,\% between Fourier and Lord--Shulman for the parameters used.

\subsection{Micro- and nano-electromechanical resonators}
Thermoelastic damping is a leading loss mechanism in flexural micro-resonators. The classical Zener relaxation (Section~\ref{sec:ted}) captures the dependence on thickness and frequency, and the extension to beams and plates by Lifshitz and Roukes \citep{Lifshitz2000} is standard. The corrections reviewed here enter in two ways. For ultra-thin structures and very high frequencies, lagging and non-Fourier effects change the damping; DPL-based thermoelastic damping in nonlocal nanoplates \citep{Yang2021}, couple-stress micro-plates \citep{Kumar2017} and rotating nano-beams with spatiotemporal nonlocality \citep{Sur2025} are examples. And at small sizes nonlocal and gradient elasticity modify the vibration frequencies on which the damping depends. Two cautions follow from our analysis: the form of the nonlocal thermoelastic constitutive law (form A or B, Section~\ref{sec:size}) changes the predicted damping near the elastic cut-off, and the lag time enters only through $\omega\tau$, so unless the frequency is within one or two decades of $1/\tau_q$, the classical result is adequate.

\subsection{Piezoelectric devices, sensors and semiconductors}
Thermopiezoelectric coupling affects bulk and surface acoustic wave devices, energy harvesters and sensors, and photothermal effects in piezoelectric semiconductors add carrier dynamics \citep{Gupta2023,Gupta2024a,Gupta2026HeatMap}. In the quasi-static, charge-free regime the effect of the electric field on the thermoelastic waves is the exact renormalisation \eqref{eq:renorm}; designs that rely on thermal sensing via the pyroelectric effect depend on the sign convention discussed there. Interface problems (substrate--film, electrode--crystal) are the subject of the energy-ratio literature \citep{Barak2023IJP,Barak2023JTS,Gupta2023Energy,Kumar2024,Kumar2024a}, where Section~\ref{sec:interfaces} suggests the checks to apply.

\subsection{Biological tissue and heterogeneous media}
Heat transfer in biological tissue is classically described by the bioheat equation of Pennes \citep{Pennes1948}, which adds blood perfusion to the Fourier equation. Non-Fourier models have been applied to tissue, based on reports of large relaxation times in nonhomogeneous materials \citep{Kaminski1990,Mitra1995}; the interpretation of the corresponding measurements has been debated, and the DPL and fractional models have been fitted to such data \citep{vanDerMerwe2021,Yu2024}. Shock-induced nonlocal coupled thermoelasticity analysis with energy dissipation in a MEMS/NEMS beam resonator was presented by \citep{Hosseini2017}. Fractional order thermo-viscoelastic theory of biological tissue with dual phase lag heat conduction model was developed by \citep{Li2021}; bio-heat transfer analysis based on fractional derivative and memory-dependent derivative heat conduction models was performed by \citep{Zhang2021}. For water-like media $\omega_*=\rho cc_1^2/k\approx10^{13}$\,rad/s (with $\rho\approx10^3$\,kg/m$^3$, $c\approx4\times10^3$\,J/kg\,K, $c_1\approx1.5\times10^3$\,m/s, $k\approx0.6$\,W/m\,K), far above any frequency of interest, and the thermoelastic coupling is weak, so the interest in tissue is in the thermal field itself rather than in thermoelastic waves.

Effect of magnetic field and gravity on 2D fiber-reinforced medium under fractional order theory of thermoelasticity was analysed by \citep{Ma2020}.

\subsection{Geophysical and porous media}
Porous and double-porous thermoelastic models with two temperatures are applied to wave propagation in saturated and partially saturated solids \citep{Singh2020,Kumar2022,Kumar2025}; here the additional slow waves, strongly attenuated, carry fluid and heat. We note that the validity of the generalized heat laws in such media, in which the relevant time scales are controlled by pore-scale processes, is largely untested.

\begin{figure}[t!]\centering\small
% fig11_applications.tex -- Application regime map (s15_applications)
\begin{tikzpicture}[x=1.0cm,y=1.0cm,>=Stealth]
  \tikzset{
    node/.style={draw,rounded corners=1pt,fill=blue!7,align=center,font=\footnotesize,inner sep=2pt,text width=2.6cm},
    arr/.style={->,thick,shorten >=1pt,shorten <=1pt},
    lbl/.style={font=\footnotesize,text=blue!60!black},
  }
  \node[node] at (-3.6, 0.6) (laser) {Laser\\[-1pt] ultrasonics\\[-1pt] ns--ps pulses};
  \node[node] at ( 0.6, 0.6) (res) {MEMS/NEMS\\[-1pt] resonators\\[-1pt] thermoelastic\\[-1pt] damping};
  \node[node] at ( 4.8, 0.6) (piezo) {Piezoelectric\\[-1pt] devices\\[-1pt] SAW/BAW\\[-1pt] energy harvesters};
  \node[node] at (-3.6,-2.4) (bio) {Biological\\[-1pt] tissue\\[-1pt] perfusion\\[-1pt] weak coupling};
  \node[node] at ( 0.6,-2.4) (geo) {Geophysical\\[-1pt] porous\\[-1pt] double-porosity\\[-1pt] slow waves};
  \draw[arr] (laser.east) -- (res.west) node[pos=0.5,above,lbl] {\small $\\omega\\tau,\\omega/\\omega_*$};
  \draw[arr] (res.east) -- (piezo.west) node[pos=0.5,above,lbl] {};
  \draw[arr] (laser.south) -- (bio.north) node[pos=0.5,left,lbl] {};
  \draw[arr] (res.south) -- (geo.north) node[pos=0.5,right,lbl] {};
\end{tikzpicture}

\caption{Application regime map: four areas ordered by the dimensionless numbers $\omega\tau$ (lag relevance) and $\omega/\omega_*$ (coupling relevance); Section~\ref{sec:synthesis} gives the quantitative criteria.}\label{fig:applications-geo}
\end{figure}

\subsection{Summary}
In all four areas the single most informative dimensionless number is $\omega\tau$ (frequency times the relevant lag), and the second is $\omega/\omega_*$. The regime map in Section~\ref{sec:synthesis} makes this quantitative.

\begin{summarybox}[Summary]
Applications divide by regime: laser ultrasonics (classical adequate for ns pulses; lagged/fractional for fs--ps), MEMS/NEMS resonators (Zener damping with nonlocal/lag corrections near cut-off), piezoelectric devices (exact stiffening renormalisation in charge-free regime), biological tissue (weak coupling, interest in thermal field), geophysical porous media (slow diffusive waves, largely untested). The key numbers are $\omega\tau$ and $\omega/\omega_*$.
\end{summarybox}
